\section{Introduction}

Graph Neural Networks (GNNs) have achieved remarkable success by leveraging a fundamental principle: connected nodes should share information. This \textit{message-passing} paradigm~\cite{gilmer2017mpnn} underlies nearly all modern GNN architectures, from Graph Convolutional Networks (GCN)~\cite{kipf2017gcn} to Graph Attention Networks (GAT)~\cite{velivckovic2018gat}. Yet this principle rests on an implicit assumption that is rarely questioned:

\begin{quote}
\textit{``Structural neighbors are the optimal source for information aggregation.''}
\end{quote}

\textbf{But what if your neighbors are not your friends?}

\subsection{The Problem: Structure-Feature Misalignment}

In many real-world graphs---particularly financial networks---structural proximity does not imply feature similarity or label agreement. Consider three illustrative scenarios:

\textbf{Scenario 1: Fraud Detection with Camouflage.} Fraudsters on financial platforms deliberately connect to legitimate users to appear trustworthy~\cite{dou2020enhancing}. Their structural neighbors are \textit{designed} to be different from themselves---aggregating information from these neighbors actively misleads the model.

\textbf{Scenario 2: High-Degree Hub Nodes.} In supply chain networks, large enterprises connect to hundreds of diverse partners spanning multiple industries, risk profiles, and transaction patterns. Standard message-passing aggregates signals from all these heterogeneous neighbors, diluting entity-specific characteristics.

\textbf{Scenario 3: Numerically-Dominated Tasks.} When prediction primarily depends on intrinsic numerical features (transaction amounts, account ages, credit scores) rather than network position, structural aggregation may introduce noise from feature-dissimilar neighbors.

These scenarios share a common pattern: \textbf{structural proximity and feature similarity are misaligned}. Standard GNNs, by aggregating only from structural neighbors, fail to distinguish these two orthogonal sources of information.

\subsection{Existing Solutions: Partial and Fragmented}

The research community has recognized specific failure modes of standard GNNs:

\textbf{Heterophily-Aware Methods.} H2GCN~\cite{zhu2020beyond} and MixHop~\cite{abu2019mixhop} address graphs where connected nodes have \textit{different labels}. They employ multi-hop aggregation to reach same-label nodes at distance 2. However, these methods focus on label heterophily rather than feature-structure misalignment, and provide no guidance for numerically-rich feature spaces.

\textbf{Attention Mechanisms.} GAT~\cite{velivckovic2018gat} and GATv2~\cite{brody2022gatv2} learn to weight neighbors differently. However, attention is still computed over \textit{structural} neighbors only---they cannot attend to feature-similar but structurally-distant nodes.

\textbf{Feature Normalization.} Various preprocessing techniques (z-score, batch normalization) handle feature scale but do not address the fundamental question of \textit{which nodes to aggregate from}.

What is missing is a \textbf{unified framework} that:
\begin{enumerate}
    \item Formally distinguishes structural and feature similarity as independent information sources
    \item Provides theoretical analysis of when each source dominates
    \item Enables principled method selection based on graph characteristics
\end{enumerate}

\subsection{Our Contribution: Feature-Structure Disentanglement}

We introduce the \textbf{Feature-Structure Disentanglement (FSD)} framework, which provides a unified perspective on graph attention design. Our key insight is:

\begin{quote}
\textit{Structural proximity and feature similarity are orthogonal signals. Optimal GNN design requires matching the aggregation strategy to the dominant signal source.}
\end{quote}

\textbf{Contribution 1: Extended Theoretical Framework.} We formalize the distinction between structure-based and feature-based similarity through the \textbf{Feature-Structure Disentanglement (FSD)} framework. Building on the \textbf{Feature-Structure Alignment Score} ($\rho_{\text{FS}}$), we introduce \textbf{Aggregation Dilution} ($\delta_{\text{agg}}$)---a novel metric capturing degree-dependent information loss during mean aggregation. We prove that $\delta_{\text{agg}} = \mathbb{E}[d_i \cdot (1 - S_{\text{feat}})]$ provides a tighter bound on when mean aggregation fails (Theorems~\ref{thm:single_layer}--\ref{thm:concat}).

\textbf{Contribution 2: Unified View of GNN Aggregation Strategies.} We show that the choice between mean aggregation (GCN, GAT), sampling (GraphSAGE), and concatenation (H2GCN) depends critically on $\delta_{\text{agg}}$:
\begin{itemize}
    \item Low $\delta_{\text{agg}}$ ($<$ 5): Mean aggregation with NAA provides numerical attention benefits
    \item High $\delta_{\text{agg}}$ ($>$ 10): Sampling bounds dilution; concatenation preserves ego information
\end{itemize}
This explains \textit{why} certain methods excel on certain graphs (Table~\ref{tab:fsd_special_cases}).

\textbf{Contribution 3: Numerically-Aware Attention (NAA).} As an instantiation of FSD for financial graphs with low aggregation dilution, we propose NAA, which explicitly incorporates feature similarity through:
\begin{itemize}
    \item Log-scale normalization to preserve numerical magnitude across orders of magnitude
    \item Explicit numerical similarity term in attention computation
    \item Adaptive head gating for regularization on smaller graphs
\end{itemize}

\textbf{Contribution 4: Comprehensive Empirical Validation.} We validate our extended framework on \textbf{five} fraud detection benchmarks (Elliptic, Amazon, YelpChi, DGraphFin, IEEE-CIS). Key findings:
\begin{itemize}
    \item On \textbf{sparse, high-dimensional graphs with low $\delta_{\text{agg}}$} (Elliptic, Amazon): NAA significantly outperforms baselines (+25\% AUC, +0.46 F1)
    \item On \textbf{dense graphs with high $\delta_{\text{agg}}$} (YelpChi, IEEE-CIS): H2GCN/GraphSAGE win by 7\% AUC---\textbf{correctly predicted by $\delta_{\text{agg}}$ alone}, while $\rho_{\text{FS}}$ would fail
    \item Under \textbf{extreme class imbalance} (DGraphFin 1.3\%): All methods fail---data challenges dominate
\end{itemize}
The extended FSD framework achieves \textbf{5/5 correct a priori predictions}, including IEEE-CIS where the original framework would fail.

\textbf{Contribution 5: Practitioner Guidelines.} We provide a three-metric diagnostic ($\rho_{\text{FS}}$, $\delta_{\text{agg}}$, $h$) for principled GNN method selection without exhaustive experimentation.

\subsection{Scope and Honest Positioning}

We emphasize that FSD is a \textbf{diagnostic framework}, not a claim of universal superiority for any single method:
\begin{itemize}
    \item NAA excels when features are high-dimensional and informative
    \item H2GCN excels when heterophily is the dominant challenge
    \item Standard GAT suffices when structure and features are well-aligned
\end{itemize}

This honest characterization is a feature, not a limitation. By understanding \textit{when} each approach works, practitioners can make informed choices rather than defaulting to the latest architecture.

\subsection{Paper Organization}

Section~\ref{sec:related} surveys related work on graph attention, heterophily, and numerical feature processing. Section~\ref{sec:fsd} presents the FSD theoretical framework with formal definitions and analysis. Section~\ref{sec:method} describes NAA as an FSD instantiation for financial graphs. Section~\ref{sec:experiments} validates our theoretical predictions across four benchmarks. Section~\ref{sec:discussion} discusses implications, limitations, and broader applicability. Section~\ref{sec:conclusion} concludes with future directions.
